\scanpage{89}
\deleted{\unsure{On|e}}et \unsure{a}
\[
\begin{aligned}
\Omega(6)&=\mathrm{Aut}(\Pi(6))\\
\widetilde{\Omega}(6)&=\mathrm{Aut}\,\mathrm{ext}(\Pi(6))
\end{aligned}
\]
de sorte qu'on a un hom \unsure{évident|induit|induits}
\diagram{\begin{tikzcd}[ampersand replacement=\&]
\Omega.\arrow[r]\arrow[d]\&\Omega(6)\arrow[d]\& \text{et un au \unsure{niv.} des}\\
\widetilde{\Omega}\arrow[r]\&\widetilde{\Omega}(6)\&
\end{tikzcd}}
Ceci dit, l'hom composé
\diagram{\begin{tikzcd}[ampersand replacement=\&,column sep=small]
{\unsure{G|\mathfrak S|\mathcal S}}_3\cdot(\gamma\times\gamma)^{(*G_3/\pi_{01})}\arrow[r,dash]\arrow[dr]\&\Omega\arrow[r,dash]\&\Omega(6)\\
\text{se factorise en}\&G_3\cdot(\gamma_0\times\gamma_0)^{(*G_3/\pi_{01})}\arrow[ur]\&
\end{tikzcd}}
i.e. les op. ``\unsure{él|èl|él.}'' au sens profini, ou au sens discret, deviennent les $=$, dès lors qu'on opère dans $\Pi(6)$ au lieu de $\Pi$. On peut dire que si $H$ est un gpe où tt él. est d'ordre divisant $6$ (p.ex. $H\simeq A_4$) alors l'op. de $G_3\underset{1/2}{\cdot}|\unsure{E|\mathcal E|\mathbb E}|\unsure{\simeq|\xrightarrow{\sim}}\Omega^+$ sur l'ens $\unsure{\mathscr T|\mathcal T}(H)$ ou $\mathrm{Sch}(H)$ se fait par l'intermédiaire de $G_3\underset{1/2}{\cdot}|E|_0\simeq\Omega_0^+$ -- donc aussi par l'intermédiaire de $G_3\underset{1/2}{\cdot}(\gamma_0\times\gamma_0)^{(*G_3/\pi_{01})}$.
% archive sheet number: 88
